\begin{rubric}{Employment history}

\entry*[Nov. 2023 -- \textbf{Present}]%
	\textbf{Member of Scientific Staff} $\bullet$ RWTH Aachen, Chair of Applied and Computational Mathematics.

Development of kinetic theory approaches for non-equilibrium flow simulations. Development of entropy-stable DG schemes for multi-component real gas flows with thermo-chemical non-equilibrium. Development and teaching of own Master course "Particle-based Simulation Methods". Teaching of Global Exercises in Math for bachelor students in the Computational Engineering and Math study program.

\entry*[Jan. 2023 -- Sep. 2023]%
	\textbf{Scientific Consultant} $\bullet$ UT Austin, Oden Institute for Computational Engineering and Sciences / Department of Department of Aerospace Engineering and Engineering Mechanics.
 
Development of new PIC-DSMC collision algorithms. Mentoring of Master student developing PIC-DSMC methods for high-voltage gap closure modeling. 

\entry*[Nov. 2021 -- Oct. 2023]%
	\textbf{Postdoctoral Researcher} $\bullet$ German Aerospace Center (DLR), Institute of Aerodynamics and Flow Technology, Spacecraft Department.

 Uncertainty quantification of hypersonic expanding continuum flows. Development of new Direct Simulation Monte Carlo methods for plasma simulation. Development of a photon Monte Carlo radiation transport code. Development of moment method models for gases with internal degrees of freedom. Work on low-rank methods for rarefied gas dynamic applications. Supported by a personal grant of the Alexander von Humboldt Foundation.

\entry*[July 2018 -- June 2021]%
	\textbf{Postdoctoral Research Fellow} $\bullet$ UT Austin, Oden Institute for Computational Engineering and Sciences / Department of Department of Aerospace Engineering and Engineering Mechanics.

 Development of a hybrid Discrete Velocity Method/Direct Simulation Monte Carlo code. Development of new Direct Simulation Monte Carlo algorithms for modeling of ionized rarefied gas flows. Development of new kinetic models for gases with internal degrees of freedom and chemically reacting gas flows.
Research project supported by Sandia National Laboratories.


\entry*[Sept. 2015 -- Dec. 2017]%
	\textbf{Research engineer} $\bullet$ Saint-Petersburg State University, Department of Hydroaeromechanics.

 Lead developer of a C++ library aimed for kinetic theory computations and coupling with CFD codes. Implementation of state-to-state models in Direct Simulation Monte Carlo (DSMC) codes, development of simplified models for vibrational relaxation rates, and numerical modeling of reaction rates in viscous gas flows.


\entry*[Apr. 2013 -- Aug. 2015]%
	\textbf{Research engineer} $\bullet$ Saint-Petersburg State University, Department of Hydroaeromechanics.

Development of theoretical models of reaction rates in viscous gas flows and their numerical modeling.
\end{rubric}

% \begin{rubric}{Berufliche Laufbahn}

% \entry*[Nov. 2023 -- \textbf{bis jetzt}]%
% 	{Wissenschaftlicher Mitarbeiter} $\bullet$ \textbf{Lehrstuhl f\"{u}r Applied and Computational Mathematics, RWTH Aachen, Aachen, Deutschland}.

% Entwicklung und Anwendung von kinetischen Verfahren f\"{u}r Modellierung von Str\"{o}mungen mit thermischen und chemischen Nichtgleichgewicht. Entwicklung von Entropy-Stabilen Discontinuous Galerkin Methoden f\"{u}r reale Gase. 

% % Development of kinetic theory approaches for non-equilibrium flow simulations. Development of entropy-stable DG schemes for multi-component real gas flows. Teaching of Global Exercises in Math for bachelor students in the Computational Engineering and Math study program.

% \entry*[Jan. 2023 -- Sep. 2023]%
% 	{Wissenschaftlicher Berater} $\bullet$ \textbf{Oden Institute for Computational Engineering and Sciences / Department of Aerospace Engineering and Engineering Mechanics, University of Texas at Austin, Austin, TX, USA}.

%  Entwicklung von neuen PIC-DSMC Kollisionsverfahren. Betreuung einer Masterstudentin, die PIC-DSMC f\"{u}r Simulation von Gap Closure im Hochspannungsdurchbruch anwandte.
% % Development of new PIC-DSMC collision algorithms. Mentoring of Master student developing PIC-DSMC methods for high-voltage gap closure modeling. 

% \entry*[Nov. 2021 -- Oct. 2023]%
% 	{Postdoktorand} $\bullet$ \textbf{Institut f\"{u}r Aerodynamik und Str\"{o}mungstechnik, Abteilung Raumfahrtszeuge, Deutsches Zentrum f\"{u}r Luft- und Raumfahrt (DLR), G\"{o}ttingen, Deutschland}.

% Quantifizierung der Unischercheit von hoch-Enthalpie expandierenden Str\"{o}mungen. Entwicklung von neuen Verfahren f
% "{u}r Simulation von Plasmen und verd\"{u}nnten Gasen.
% Forschungsprojekt gef\"{o}rdert durch ein Stipendium der Alexander von Humboldt Stiftung.
%  % Uncertainty quantification of hypersonic expanding continuum flows. Development of new Direct Simulation Monte Carlo methods for plasma simulation. Development of a photon Monte Carlo radiation transport code. Development of moment method models for gases with internal degrees of freedom. Work on low-rank methods for rarefied gas dynamic applications. Supported by a personal grant of the Alexander von Humboldt Foundation.

% \entry*[Jul. 2018 -- Jun. 2021]%
% 	{Postdoktorand} $\bullet$ \textbf{Oden Institute for Computational Engineering and Sciences / Department of Aerospace Engineering and Engineering Mechanics, University of Texas at Austin, Austin, TX, USA}.

% Entwicklung eines hybriden Discrete Velocity Method/Direct Simulation Monte Carlo Verfahrens. Entwicklung von neuen Direct Simulation Monte Carlo Algorithmen f\"{u}r Modellierung von verd\"{u}nnten ionisierten  Gasstr\"{o}mungen.
% %  Development of a hybrid Discrete Velocity Method/Direct Simulation Monte Carlo code. Development of new Direct Simulation Monte Carlo algorithms for modeling of ionized rarefied gas flows. Development of new kinetic models for gases with internal degrees of freedom and chemically reacting gas flows.
% % Research project supported by Sandia National Laboratories.

% \entry*[Nov. 2017 -- Jan. 2018]%
% 	{Forschungspraktikum} $\bullet$ \textbf{Institut f\"{u}r Aerodynamik und Str\"{o}mungstechnik, Abteilung Raumfahrtszeuge, Deutsches Zentrum f\"{u}r Luft- und Raumfahrt (DLR), G\"{o}ttingen, Deutschland}.
%     \par Umzetungs eines neuen Models f\"{u}r Vibrationsrelaxation in Luftmischungen in dem DLR-TAU CFD L\"{o}ser. Forschungsprojekt gef\"{o}rdert durch ein DAAD-Stipendium.
    
% \entry*[Oct. 2016]%
% 	{Forschungspraktikum} $\bullet$ \textbf{Bundesuniversität von Paraná, Curitiba, Brazil}.
%     \par Forschung von Wellenausbreitung in Study of wave propagation processes in Hochtemperatur-Gasstr\"{o}mungen. Forschungsprojekt gef\"{o}rdert durch ein Stipendium der Santander Bank.
    
% \entry*[Feb. 2016]%
% 	Forschungspraktikum $\bullet$ \textbf{Khristianovich Institut f\"{u}r Theoretische und Angewandte Mechanik, Novosibirsk, Russland}.
%     \par Umsetzung von State-to-state Modellen von inelastischen Prozessen in einem DSMC Code. Forschungsprojekt gef\"{o}rdert durch ein Stipendium des Russischen Fonds f\"{u}r Grundlagenforschung.
    
% % \entry*[Sept. 2015 -- Dec. 2017]%
% \entry*[Apr. 2013 -- Dec. 2017]%
% 	{Wissenschaftlicher Mitarbeiter} $\bullet$ \textbf{Lehrstuhl f\"{u}r Hydroaeromechanik, Staatliche Universit\"{a}t Sankt-Petersburg, Sankt-Petersburg, Russland}. 

%  Entwicklung von theoretischen Modellen von Reaktionsraten in viskosen Gasstr\"{o}mungen und ihre numerische Modellierung.
%  Hauptentwickler von einer C++ Bibliothek f\"{u}r Berechnungen von kinetischen Transportkoeffizient und Reaktionsraten, gezielt auf Kopplung mit CFD (Computational Fluid Dynamics) Codes.
% %  Lead developer of a C++ library aimed for kinetic theory computations and coupling with CFD codes. 
% % Development of theoretical models of reaction rates in viscous gas flows and their numerical modeling. Implementation of state-to-state models in Direct Simulation Monte Carlo (DSMC) codes, development of simplified models for vibrational relaxation rates, and numerical modeling of reaction rates in viscous gas flows.


% % \entry*[Apr. 2013 -- Aug. 2015]%
% % 	\textbf{Research engineer} $\bullet$ Saint-Petersburg State University, Department of Hydroaeromechanics.

% \end{rubric}